%==============================================================================
% IMPORTANT: Every numerical table and figure in this file must be regenerated
% after changing sensing_min_power_fraction. Do not use legacy zero-lower-bound
% results as evidence for the participation-constrained model.
\section{Numerical Results}\label{sec:numerical}
%==============================================================================

\subsection{Simulation Setup}
Unless otherwise stated, the Cell-Free ISAC network contains $M=6$ APs with
$N_t=2$ antennas per AP, $K=3$ UEs, and $P=2$ targets. Each target state has
dimension $N_\theta=2$, the deployment area is $400\,\mathrm{m}\times
400\,\mathrm{m}$, and the AP power limit is $20$ dBm ($0.1$ W). The nominal
CSI uncertainty radius is $\epsilon_h=0.05$, the communication and sensing
SINR thresholds are $0$ dB, and the automatically calibrated PCRB allowance
uses $\Gamma_{\mathrm{alpha}}=3$. A result is counted as feasible only after
binary-topology recovery, fixed-$b$ re-optimization, and physical validation.

\begin{table}[htbp]
\centering
\caption{Main simulation configuration.}
\label{tab:simulation-setup}
\begin{tabular}{lc}
\toprule
Parameter & Value\\
\midrule
APs and antennas per AP & $M=6$, $N_t=2$\\
UEs, targets, and state dimension & $K=3$, $P=2$, $N_\theta=2$\\
Required sensing APs per target & $N_{\mathrm{req}}=3$\\
Deployment area and AP power limit & $400\,\mathrm{m}\times400\,\mathrm{m}$, $0.1$ W\\
CSI uncertainty radius & $\epsilon_h=0.05$\\
Communication and sensing SINR targets & $0$ dB, $0$ dB\\
PCRB allowance & Automatic calibration, $\Gamma_{\mathrm{alpha}}=3$\\
Dedicated sensing participation floor & $P_{\min}^{\rm sen}=0.01P_{\max}=1$ mW\\
Solver and implementation & MATLAB R2024b, CVX, MOSEK 11.2.2\\
\bottomrule
\end{tabular}
\end{table}

\subsection{Binary-DC Recovery and Topology Quality}
The binary penalty is evaluated over $30$ common scenarios. Relative to the
unpenalized continuous relaxation, the median binary residual decreases from
$4.268\times10^{-1}$ to $6.229\times10^{-5}$ under the dual-DC update.
In this MU--MISO configuration, the transmit covariance is already nearly
rank one after SDR; consequently, the rank penalty is retained as a general
formulation component but is not claimed to yield an additional empirical
gain in this setting. Every reported physical point is obtained only after
binary topology recovery, fixed-$b$ re-optimization, and explicit validation.

\begin{table}[htbp]
\centering
\caption{Current-model topology and recovery evidence over $30$ common scenarios.}
\label{tab:main-mc}
\begin{tabular}{lc}
\toprule
Metric & Result\\
\midrule
Binary residual, unpenalized / dual-DC median & $4.268\times10^{-1}$ / $6.229\times10^{-5}$\\
FIM fixed-topology feasibility & $30/30$\\
FIM versus nearest-AP mean power reduction ($N_{\rm req}=3$) & $21.27\%$\\
FIM versus random: jointly feasible cases ($N_{\rm req}=3$) & $10/10$; mean reduction $73.30\%$\\
FIM power gap to SDR, median & $9.73\%$\\
\bottomrule
\end{tabular}
\end{table}

Fig.~\ref{fig:architecture} depicts the asymmetric signaling architecture.
The communication covariances $\{\mat{W}_k\}$ remain globally cooperative,
whereas $b_{mp}$ requires each selected AP to allocate nonzero dedicated
sensing covariance $\mat{S}_p$ to target $p$. The topology comparison shows
that the FIM-aware candidate is not interchangeable with a path-loss-only or
random association: at the main setting $N_{\rm req}=3$, its mean power is
$21.27\%$ below the nearest-AP baseline, while random association is feasible
in only $10$ of $30$ trials. Support
stability is used only as an engineering stopping rule; final binary
feasibility is certified by fixed-$b$ re-optimization.

\begin{figure}[htbp]
\centering
\includegraphics[width=\linewidth]{../figures/fig1_system_architecture.png}
\caption{Cell-Free ISAC architecture. Global communication beamforming is not
gated by $b_{mp}$, whereas each selected AP--target pair radiates a
nonzero target-specific dedicated sensing component.}
\label{fig:architecture}
\end{figure}

\subsection{Cluster Size, QoS, and Recovery}
Fig.~\ref{fig:cluster-size} uses $30$ common seeds for every
$N_{\mathrm{req}}\in\{2,3,4,5,6\}$. The mean power decreases from
$37.65$ mW at $N_{\mathrm{req}}=2$ to $29.50$ mW at
$N_{\mathrm{req}}=4$, then rises to $30.94$ mW at $N_{\mathrm{req}}=6$.
Thus, the extra geometric diversity initially reduces tracking cost, whereas
the mandatory per-pair sensing floor eventually offsets this gain. All tested
points are physically feasible. The QoS
statistics in Fig.~\ref{fig:qos-cluster} show that the PCRB constraint is
numerically active: the maximum target PCRB ratio is essentially one. The
worst-user nominal communication and sensing-SINR margins remain
nonnegative. A recomputation from the returned covariance matrices gives a
maximum PCRB ratio numerically equal to one at the displayed precision.

\begin{figure}[htbp]
\centering
\includegraphics[width=\linewidth]{../figures/fig3_cluster_size_tradeoff.png}
\caption{Physical feasibility, total transmit power, and runtime versus
the required sensing-cluster size. Error bars show the 10th--90th
percentile interval.}
\label{fig:cluster-size}
\end{figure}

\begin{figure}[htbp]
\centering
\includegraphics[width=\linewidth]{../figures/fig4_qos_vs_cluster_size.png}
\caption{QoS statistics versus sensing-cluster size. Shaded regions denote
the 10th--90th percentile interval across feasible trials.}
\label{fig:qos-cluster}
\end{figure}

\subsection{Comparison with Alternative Association Rules}
Fig.~\ref{fig:method-comparison} compares the proposed binary-DC recovery
with an FIM-greedy topology, a nearest-AP topology, and a random cardinality-
feasible topology over $30$ common scenarios for each value of
$N_{\mathrm{req}}$. After a topology is selected, all physical baselines use
the same fixed-assignment robust beamforming and dedicated sensing-covariance
optimization. Hence, the comparison isolates the association and recovery
mechanism rather than attributing a continuous beamforming advantage to a
discrete rule. The SDR is retained only as a power lower bound and is not
included as a physical curve.

The proposed method is physically feasible in all $30$ trials at every
$N_{\mathrm{req}}$. The FIM-greedy rule is also feasible in all trials, but
its mean power is slightly higher. In contrast, the nearest-AP rule is
feasible in only $24/30$ trials at $N_{\mathrm{req}}=2$, while the random
rule is feasible in only $1/30$ trials. At $N_{\mathrm{req}}=3$, the proposed
method uses $30.83$ mW on average, compared with $30.93$ mW for FIM-greedy
and $39.29$ mW for nearest-AP association. The random baseline must be read
conditional on feasibility: it is feasible in only $10/30$ trials and has a
conditional mean power of $115.82$ mW. The average normalized PCRB remains
essentially one for every feasible method, confirming that the power gains do
not result from relaxed tracking quality. The reported sum rate is the
nominal aggregate spectral efficiency of the returned physical covariance;
it is descriptive rather than an additional optimization objective.

For every feasible physical solution, the number of nonzero AP--target
dedicated-sensing pairs equals $PN_{\mathrm{req}}$, as required by the
positive participation floor. The final panel of Fig.~\ref{fig:method-comparison}
therefore reports the number of \emph{unique} APs with nonzero sensing power,
which captures AP reuse across targets. At $N_{\mathrm{req}}=6$, every AP is
assigned to every target and all association rules reduce to the same topology
and the same solution.

\begin{figure*}[htbp]
\centering
\includegraphics[width=\textwidth]{../figures/fig10_method_comparison_vs_nreq.png}
\caption{Thirty-seed comparison of physical association and recovery methods
versus the required sensing-cluster size. Power, sum rate, sensing SINR, and
the number of active sensing APs are conditional on physical feasibility;
shaded bands denote the 10th--90th percentile interval.}
\label{fig:method-comparison}
\end{figure*}

For the same $30$ scenarios, the complete recovery pool achieves a mean power
of $30.83$ mW, compared with $30.93$ mW for the FIM fixed topology. The small
difference indicates that FIM geometry is already an effective topology
selector in the present architecture; the continuous robust beamforming and
dedicated sensing-covariance optimization then supply the remaining gain.
Fig.~\ref{fig:dimension} confirms the expected runtime increase with total
antenna dimension.

\begin{figure}[htbp]
\centering
\includegraphics[width=.78\linewidth]{../figures/fig7_dimension_sensitivity.png}
\caption{Runtime and physical feasibility versus total antenna dimension.}
\label{fig:dimension}
\end{figure}

\subsection{Communication--Sensing Trade-Off and CSI Robustness}
Fig.~\ref{fig:tradeoff-mc} averages a $5\times5$ QoS grid over five common
seeds. Mean transmit power increases when the communication SINR requirement
is tightened or the PCRB allowance scale $\alpha$ is reduced. All $25$ tested
grid points are physically feasible; hence the heat map should be interpreted
as evidence of feasibility over this tested range, rather than as a boundary
of the full feasible region.

Fig.~\ref{fig:robustness} compares the proposed robust design with a nominal
CSI baseline that removes the S-procedure and is designed with
$\epsilon_h=0$. Both designs are tested under identical independently sampled
complex channel perturbations inside the uncertainty ball. For
$\epsilon_h\in\{0.02,0.05,0.08\}$, the robust design has zero empirical
system outage in the $100$ perturbation samples per seed over $30$ seeds,
whereas the mean nominal-design outage is $89.0\%$, $89.6\%$, and $91.9\%$,
respectively. At the uncertainty radius $0.08$, the median robust power premium is only
$5.12\%$.
This sampled result complements, but does not replace, the S-procedure
certificate for the prescribed uncertainty set.

Finally, a $30$-seed participation-floor sensitivity study verifies the
physical interpretation of $b_{mp}$. For each positive floor in
$P_{\min}^{\rm sen}/P_{\max}\in\{0.5\%,1\%,2\%,5\%\}$, all $30$ returned
solutions satisfy the corresponding per-AP--target lower bound with zero
measured violation. The mean power is $30.83$ mW at the selected $1\%$ floor,
whereas a $5\%$ floor raises it to $37.11$ mW. This supports $1\%$ as a
non-degenerate participation level with modest energy cost.

\begin{figure}[htbp]
\centering
\includegraphics[width=\linewidth]{../figures/fig8_statistical_tradeoff.png}
\caption{Statistical communication--sensing trade-off surface and physical
feasibility heat map over five common seeds.}
\label{fig:tradeoff-mc}
\end{figure}

\begin{figure}[htbp]
\centering
\includegraphics[width=\linewidth]{../figures/fig9_csi_robustness.png}
\caption{Robust and nominal CSI designs under ball-bounded channel
perturbations.}
\label{fig:robustness}
\end{figure}
